{{دیوان
| عنوان = ۔۔۔۔ کی گود میں بِلّی دیکھ کر
| ماخذ = ویکی ماخذ
| ماخذ_لنک = https://ur.wikisource.org/wiki/%DA%A9%DB%8C_%DA%AF%D9%88%D8%AF_%D9%85%DB%8C%DA%BA_%D8%A8%D9%90%D9%84%D9%91%DB%8C_%D8%AF%DB%8C%DA%A9%DA%BE_%DA%A9%D8%B1
| مقابلہ = اقبال سائبر لائبریری (اقبال اکادمی پاکستان)
| مقابلہ_لنک = https://www.iqbalcyberlibrary.net
}}
<poem>
تجھ کو دُزدیدہ نگاہی یہ سِکھا دی کس نے
رمز آغازِ محبّت کی بتا دی کس نے

ہر ادا سے تری پیدا ہے محبّت کیسی
نیلی آنکھوں سے ٹپکتی ہے ذکاوت کیسی

دیکھتی ہے کبھی ان کو، کبھی شرماتی ہے
کبھی اُٹھتی ہے، کبھی لیٹ کے سو جاتی ہے

آنکھ تیری صفَتِ آئنہ حیران ہے کیا
نُورِ آگاہی سے روشن تری پہچان ہے کیا

مارتی ہے انھیں پونہچوں سے، عجب ناز ہے یہ
چھیڑ ہے، غُصّہ ہے یا پیار کا انداز ہے یہ؟

شوخ تُو ہوگی تو گودی سے اُتاریں گے تجھے
گِر گیا پھُول جو سینے کا تو ماریں گے تجھے

کیا تجسّس ہے تجھے، کس کی تمنّائی ہے
آہ! کیا تُو بھی اسی چیز کی سودائی ہے

خاص انسان سے کچھ حُسن کا احساس نہیں
صُورتِ دل ہے یہ ہر چیز کے باطن میں مکیں

شیشۂ دہر میں مانندِ مئے ناب ہے عشق
رُوح خورشید ہے، خونِ رگِ مہتاب ہے عشق

دلِ ہر ذرّہ میں پوشیدہ کسک ہے اس کی
نور یہ وہ ہے کہ ہر شے میں جھلک ہے اس کی

کہیں سامانِ مسرّت، کہیں سازِ غم ہے
کہیں گوہر ہے، کہیں اشک، کہیں شبنم ہے
</poem>
